\subsection{Modules over a Semisimple Ring}

PROPOSITION 4. --- Let A be a ring. The following properties are equivalent:

\begin{enumerate}
\def\labelenumi{(\roman{enumi})}
\item
  The ring A is semisimple.
\item
  Every A-module is semisimple.
\item
  There exists a generating and semisimple A-module.
\item
  There exists a faithful and semisimple A-module with finitely generated countermodule.
\item
  Every A-module is projective.
\item
  Every monogenous A-module is projective.
\end{enumerate}

*For other characterizations of semisimple rings, see Proposition 6 of X, §8, n° 4, p.~140.*

Let us first prove that (i) implies (ii) and (v). Suppose that the ring A is semisimple, and consider a left A-module M. By assumption, the A-module \(A_{s}\) is semisimple, so every free A-module is semisimple. By Proposition 20 of II, §1, No.~11, p.~218, there exist a free A-module L and a surjective A-linear mapping u from L to M. Let N be the kernel of u. Since the A-module L is semisimple, there exists a semisimple submodule \(N'\) supplementary to N in L (VIII, p.~56, Theorem 1). The A-module \(N'\) is projective, and u induces an isomorphism from \(N'\) to M. Consequently, M is semisimple and projective.

\begin{enumerate}
\def\labelenumi{(\roman{enumi})}
\setcounter{enumi}{1}
\item
  \(\Rightarrow\) (iii): If every A-module is semisimple, then the A-module \(A_{s}\) is semisimple; it is moreover generating.
\item
  \(\Rightarrow\) (iv): If M is a generating A-module, then it is faithful (VIII, p.~80, Corollary of Theorem 1), and its countermodule is finitely generated (VIII, p.~99, Corollary 1).
\item
  \(\Rightarrow\) (i): Let M be a faithful and semisimple A-module with finitely generated countermodule. By Lemma 4 of VIII, p.~8, there exists a natural number m such that \(A_{s}\) is isomorphic to a submodule of \(M^{m}\) . The A-module \(M^{m}\) is semisimple, and therefore so is \(A_{s}\) .
\end{enumerate}

The implication \((v) \Rightarrow (vi)\) is immediate.

\begin{enumerate}
\def\labelenumi{(\roman{enumi})}
\setcounter{enumi}{5}
\tightlist
\item
  \(\Rightarrow\) (i): Suppose that every monogenous A-module is projective. Let \(\mathfrak{a}\) be a left ideal of A. Since the A-module \(A_s / \mathfrak{a}\) is projective, there exists a left ideal \(\mathfrak{b}\) of A such that \(A_s\) is the direct sum of \(\mathfrak{a}\) and \(\mathfrak{b}\) (II, §2, No.~2, p.~231, Proposition 4). Consequently, the ring A is semisimple (VIII, p.~135, Theorem 1).
\end{enumerate}

Lemma 1. --- Let A be a left Artinian ring, and let M be a simple A-module. Then the ring \(A_{M}\) is simple.

Since the ring A is left Artinian, the same holds for the ring \(A_{M}\) , by Proposition 5 of VIII, p.~7. Now, M is a faithful and simple \(A_{M}\) -module, so the ring \(A_{M}\) is simple (VIII, p.~119, Proposition 1).

PROPOSITION 5. --- Suppose that the ring A is semisimple. Let M be a left A-module. The countermodule of M is finitely generated, and we have \(A_{M} = A_{M}^{\prime\prime}\) .

First consider the case when M is a simple A-module. By Lemma 1, the ring \(A_{M}\) is simple. Proposition 5 then follows from Lemma 1 of VIII, p.~120.

Now consider the general case. The set S of classes of simple A-modules is finite (VIII, p.~136, Proposition 1). For every \(\lambda \in S\) , choose an A-module \(S_{\lambda}\) of class \(\lambda\) , and denote the opposite field of the commutant of \(S_{\lambda}\) by \(D_{\lambda}\) . By Lemma 1 applied to the simple ring \(A_{S_{\lambda}}\) , \(S_{\lambda}\) is a finite-dimensional vector space over the field \(D_{\lambda}\) ; denote its dimension by \(m(\lambda)\) . Let B be the opposite ring of the endomorphism ring of M. We have seen in VIII, p.~84 that there exist \((D_{\lambda}, B)\) -bimodules \(V_{\lambda}\) that are simple as B-modules and an isomorphism of \((A, B)\) -bimodules from M to \(\oplus_{\lambda \in S} S_{\lambda} \otimes_{D_{\lambda}} V_{\lambda}\) . As a B-module, M is isomorphic to \(\oplus_{\lambda \in S} V_{\lambda}^{m(\lambda)}\) . Since S and the \(m(\lambda)\) are finite, M is a finitely generated B-module.

The A-module M is semisimple, and its countermodule is finitely generated. We therefore have \(A_{M} = A_{M}^{\prime\prime}\) by Proposition 4 of VIII, p.~83.

PROPOSITION 6. --- Let M be a finitely generated semisimple A-module. Denote its support (VIII, p.~66) by \(\mathcal{S}_{\mathrm{M}}\) and its endomorphism ring by B. For every \(\lambda \in \mathcal{S}_{\mathrm{M}}\) , choose a simple A-module \(S_{\lambda}\) of class \(\lambda\) , and denote the left B-module \(\mathrm{Hom}_{\mathrm{A}}(S_{\lambda}, \mathrm{M})\) by \(V_{\lambda}\) .

\begin{enumerate}
\def\labelenumi{\alph{enumi})}
\item
  The ring B is semisimple.
\item
  The mapping \(\lambda \mapsto \operatorname{cl}(V_{\lambda})\) is a bijection from the support \(\mathcal{S}_{\mathrm{M}}\) of M to the set of classes of simple B-modules.
\item
  For every \(\lambda \in \mathcal{S}_{\mathrm{M}}\) , the isotypical component of type \(\lambda\) of the A-module M is equal to the isotypical component of type \(V_{\lambda}\) of the B-module M.
\end{enumerate}

Viewed as a B-module, M is semisimple (VIII, p.~85, Proposition 5) and faithful. Its countermodule is finitely generated because we have \(A_{M} \subset \operatorname{End}_{B}(M)\) . Hence, the ring B is semisimple (Proposition 4). If \((x_{1}, \ldots, x_{r})\) is a generating sequence of the A-module M, then the mapping \(b \mapsto (bx_{1}, \ldots, bx_{r})\) from \(B_{s}\) to \(M^{r}\) is B-linear and injective. Every simple B-module is isomorphic to a submodule of \(B_{s}\) (VIII, p.~136, Proposition 1), hence to a B-submodule of M. The proposition then follows immediately from Proposition 5 of VIII, p.~85.

PROPOSITION 7. --- Let A be a semisimple ring.

\begin{enumerate}
\def\labelenumi{\alph{enumi})}
\item
  Every finitely generated A-module is reflexive (II, §2, No.~7, p.~239).
\item
  For every simple left A-module S, the right A-module \(S^{*}\) dual to S is simple, and the mapping \(\lambda \mapsto \operatorname{cl}(\lambda^{*})\) defines a bijection from the set of classes of simple A-modules to the set of classes of simple right A-modules.
\item
  Let M be a finitely generated left A-module. The right A-module \(M^{*}\) dual to M is finitely generated and has the same length as M. Moreover, we have \([M:S]=[M^{*}:S^{*}]\) for every simple A-module S.
\end{enumerate}

Let \(\mathrm{M}\) be a finitely generated A-module.

By Proposition 4 of VIII, p.~138, the A-module M is finitely generated and projective; it is therefore reflexive by Corollary 4 of II, §2, No.~7, p.~240. In particular, every simple A-module is reflexive. It also follows that two finitely generated left modules are isomorphic if and only if their duals are.

Let S be a simple left A-module. Let \((\mathrm{T}_{i})_{i\in\mathrm{I}}\) be a family of simple right A-modules with direct sum the dual \(S^{*}\) of S. Since S is reflexive, it is isomorphic to \((\mathrm{S}^{*})^{*}\) and therefore to \(\prod_{i\in I}T_{i}^{*}\) . Each of the modules \(T_{i}\) is reflexive; in particular, we have \(T_{i}^{*}\neq0\) for every \(i\in I\) . Since the simple module S is isomorphic to \(\prod_{i\in I}T_{i}^{*}\) , the set I has a single element, so \(S^{*}\) is simple.

Since M is semisimple and finitely generated, it is the direct sum of simple submodules \(S_{1},\ldots,S_{r}\) . Then \(M^{*}\) is isomorphic to the direct sum of the family \(S_{1}^{*},\ldots,S_{r}^{*}\) , and we have just seen that the modules \(S_{i}^{*}\) are simple. Assertion c) follows immediately.

